\hmsubsection{Robotics}{UMA}{OAH}\label{sec:robotics}
Robotics is a multidisciplinary field concerned with the design 
and operation of systems that interact with their physical 
environment through coordinated sensing, computation, and 
actuation \cite{craig2005}. A robotic system typically operates 
on a continuous sense--process--act cycle: the system perceives 
its surroundings through sensors, processes the resulting data to 
form a decision, and produces a physical response through 
actuators. An autonomous robotic system performs this cycle 
without direct human intervention, relying on internal logic 
to interpret sensor data and select an appropriate action.

The system developed in this project follows this principle. An 
OAK-D~S2 RGB camera mounted on the wrist and positioned above the 
workspace when the arm is in the scan pose acts as the sensor, 
observing the workspace and the objects placed within it. The host 
processor analyses the resulting image data to identify the colour 
and position of each object. The robotic arm then translates this 
decision into motion, transporting each object to the bin that 
corresponds to its detected colour. The cycle repeats continuously 
while objects remain in the workspace.

A complete robotic system therefore depends on the integration of 
several disciplines, including computer vision, embedded computing, 
mechanical design, and control. Subsequent sections describe the 
specific aspects of each discipline that are relevant to the 
present project.
